\chapter{数据口径与估值资格门槛}

\section{时间口径}

2026年8月11日09:22的采集处于盘前，腾讯与新浪行情均显示零成交。因此所有估值、涨跌幅、成交额和价格位置均使用最新完整交易日，即2026年8月10日收盘。盘前指示价单独保存，只用于证明没有把8月11日误写为收盘价。

\section{证据层级}

市场数据以腾讯批量行情、复权日线与搜狐日线成交额交叉核验。财务数据只使用CNINFO查询结果与发行人PDF。券商资料按原始PDF、券商页面/可核验摘要、目录行和未找到四级处理；只有原始PDF中的显式目标可作为外部比较锚，且其权重不等同于投资组合权重。

\section{四道门槛}

\begin{exhibitbox}[图表2｜从历史空间到当前研究结论的四道门槛]
\begin{enumerate}
  \item \textbf{价格门槛：} 当前收盘价仍应处于可解释的风险回报位置，不能把急涨后的价格称为低位。
  \item \textbf{分母门槛：} 收入、归母利润、扣非、并表与调整后利润需要同口径、同期间或明确调节。
  \item \textbf{现金门槛：} 经营现金流与营运资本必须能解释利润；负现金流并非自动否定利润，但会阻断无条件重估。
  \item \textbf{锚门槛：} 需要当前、可审计、方法可见的外部锚或可复建的自有模型；没有锚不以历史情景价替代。
\end{enumerate}
\sourcenote{门槛定义见数据审计与估值资格文件。}
\end{exhibitbox}

只要一项未通过，正确产出就是事件验证或观察，而不是给出精确目标价。本期四只原主标的均有至少一项未通过，故没有新的AStock价格区间。
